% !TeX program = xelatex
% ============================================================
% interim_report/interim_report.tex (EN) — "Thesis report" — the CP2
% deliverable of a research (academic) thesis, SE programme.
%
% Genre per the "Thesis CP2 preparation instruction 2025/26": a GOST
% 7.32-2017 document, AT LEAST 15–20 pages including the reference list
% (>= 5 GOST-formatted sources). Carries the elaborated research
% Approach. Sections mirror the instruction's list exactly; a part with
% no information is omitted while the ORDER of the rest is preserved.
% The title page is signed by the student and the supervisor.
%
% This is NOT the thesis text (that is the `thesis` document). The
% thesis report extends the CP1 Annotation with the experimental
% base, the data description and the experiment approach.
% ============================================================
\documentclass[12pt,a4paper]{extarticle}
% !TeX program = xelatex
% ============================================================
% common/annotation_preamble.tex — преамбула Аннотации (документ КТ1
% исследовательской ВКР ОП «Программная инженерия»).
%
% ЭТО НЕ ЕСПД-документ: без штампа/рамки, без кода документа, без листа
% утверждения и без содержания. Оформление по ГОСТ 7.32-2017: поля
% 30/15/20/20, Times New Roman 12 пт, интервал 1,5, обычный номер
% страницы снизу. Данные — из common/meta (\hse-макросы). Ссылки [n]
% печатаются В СТРОКУ, список источников — по ГОСТ Р 7.0.5-2008.
% ============================================================

% ── Шрифты (автогенерируются из настроек проекта) ───────────
\newcommand{\hseDefaultMono}{Courier New}
\input{../common/fonts}
\usepackage[english,russian]{babel}

% ── Геометрия по ГОСТ 7.32-2017 §6.1.1 (без места под штамп) ─
\usepackage[a4paper,left=30mm,right=15mm,top=20mm,bottom=20mm]{geometry}
\usepackage{setspace}
\onehalfspacing
\usepackage{indentfirst}
\setlength{\parindent}{1.25cm}
\setlength{\parskip}{6pt}

% ── Типографика и таблицы ───────────────────────────────────
\usepackage{microtype}
\usepackage{csquotes}
\usepackage{enumitem}
\setlist[itemize,1]{label=--}
\usepackage{graphicx}
\usepackage[table]{xcolor}
\usepackage{array}
\usepackage{tabularx}
\usepackage{booktabs}
\usepackage{float}

% ── URL и переносы ──────────────────────────────────────────
\usepackage{url}
\usepackage{xurl}
\urlstyle{same}
\usepackage{hyphenat}

% ── Цитирование: ссылки [n] В СТРОКУ (как в реальных аннотациях),
%    список источников — ГОСТ Р 7.0.5-2008. НЕ переопределяем \@cite
%    в надстрочный вид (в отличие от ЕСПД-преамбулы). ────────────
\usepackage{cite}
\makeatletter
\renewcommand\@biblabel[1]{#1.}
\makeatother

\usepackage{lastpage}
\usepackage{hyperref}
\hypersetup{
    unicode=true,
    breaklinks=true,
    colorlinks=true,
    linkcolor=black,
    urlcolor=blue,
    citecolor=blue,
    pdfborder={0 0 0}
}

% ── Данные проекта (\hse…), затем «человеческие» имена (commands) ──
% meta даёт \hseFill/\hseOptional и все \hse-макросы; commands —
% \signdateblank, \hseExecutorsBlock-зависимости (\authorsignatureline,
% \studentgroup, \studentname) для блока «Выполнил» в титуле.
\input{../common/meta}
\input{../common/commands}

% Без штампа: обычный номер страницы снизу по центру (ГОСТ 7.32 §6.3.1).
\pagestyle{plain}

\begin{document}
\sloppy
\input{title}


\begin{center}
  \textbf{\large Thesis Report}
\end{center}

\vspace{0.4em}

% ── Keywords (GOST 7.32-2017): 5–15 words or phrases, UPPERCASE,
%    nominative case, comma-separated, no full stop at the end. ──
\noindent\textbf{Keywords:} \hseFill{5–15 WORDS IN UPPERCASE, NOMINATIVE CASE, SEPARATED BY COMMAS, NO FULL STOP AT THE END}

\section{Assessment of the current state of the problem}
The problem under consideration is~\hseFill{state the scientific problem}. Today it is addressed with \hseFill{existing approaches, methods, works with citations} \cite{example-en}. Their main limitations are: \hseFill{limitation 1}; \hseFill{limitation 2}. \hseFill{Expand the review: comparison of approaches, tables, conclusions — this is the main survey section of the report}.

\section{Basis and initial data for the research}
The work is based on the curriculum of the «\hseSpecialization» programme and the thesis topic approved by the HSE FCS order. Initial data: \hseFill{existing groundwork, data, publications and software the work builds upon}.

\section{Relevance and novelty of the research topic}
The relevance stems from \hseFill{the practical and scientific significance of the task}. The scientific novelty of the work lies in~\hseFill{what exactly is proposed for the first time}.

\section{Object and subject of research}
\textbf{Object of research (or development)}~--- \hseFill{what is studied as a whole}.

\textbf{Subject of research}~--- \hseFill{the specific aspect of the object on which the work focuses}.

\section{Aim of the research}
\hseFill{a single formulation of the scientific result to be achieved}.

\section{Research objectives}
\begin{enumerate}[label=\arabic*), leftmargin=1.3cm, topsep=2pt, itemsep=2pt]
  \item \hseFill{carry out a review of existing approaches};
  \item \hseFill{propose and formalise a method (model, algorithm)};
  \item \hseFill{implement a software tool for the experiments};
  \item \hseFill{run a computational experiment and evaluate the results};
  \item \hseFill{draw conclusions on the applicability}.
\end{enumerate}

\section{Methods and methodology of the research}
\hseFill{describe the methods and the overall methodology: analytical review, formal modelling, computational experiment, statistical analysis; how they map to the objectives}.

\section{Expected new results of the research}
\begin{enumerate}[label=\arabic*), leftmargin=1.3cm, topsep=2pt, itemsep=2pt]
  \item \hseFill{method / model / algorithm and its novelty};
  \item \hseFill{software implementation for running the experiments};
  \item \hseFill{experimental evaluation and comparison with baselines};
  \item \hseFill{conclusions on the limits of applicability}.
\end{enumerate}

% ── The elaborated Approach (the core of CP2): base, data, experiments ──
\section{Experimental base of the research}
\hseFill{what the research runs on: computational environment, hardware, software infrastructure, testbeds}.

\section{Description of the research data}
\hseFill{the data the research will be run on: source, volume, structure, acquisition and preparation, usage constraints}.

\section{Approach to conducting experiments and interpreting results}
\hseFill{the experiment plan: which series, which quality metrics, baseline methods for comparison, how the results will be interpreted and validated (statistical significance, reproducibility)}.

\section{Field of application of the research results}
\hseFill{specify the field of application}.

\section{Expected adoption of the research results}
\hseFill{where and in what form the results are expected to be adopted: a product, a service, a publication, open-source code}.

\section{Economic efficiency or significance of the research}
\hseFill{where and by whom the results are applicable; what the practical or economic gain is}.

\section{Forecast assumptions on the development of the object of research}
\hseFill{how the object of research may evolve and how the work may be continued}.

\bigskip
% ── References in IEEE style (§2.5.3.1: English research theses cite in
%    IEEE style): numbered [n], author initials first, article titles in
%    quotes, venue/volume/number/pages, accessed date for online
%    sources. Every source must be cited with \cite{…} in the text.
%    CP2 instruction: AT LEAST FIVE sources. ──
\noindent\textbf{References}
\vspace{-0.3em}
\renewcommand{\refname}{}
\begin{thebibliography}{99}
  \bibitem{example-en}
  A.~B.~Author and C.~D.~Coauthor, \enquote{Title of the paper,} in \emph{Proc. of the Conference}, New York, NY, USA: ACM, 2024, pp.~10--20.

  \bibitem{example-web}
  Resource Name. (2026). Accessed: Jan.~1, 2026. [Online]. Available: \url{https://example.org}
\end{thebibliography}


\end{document}
